Este trabajo estudia la detección automática de crisis epilépticas en el electroencefalograma (EEG) con dos bases
públicas: Bonn (500 segmentos de un canal) y CHB-MIT (6 pacientes, 223 horas de registro continuo de 18 canales y 34
crisis, que ocupan el 0,37~\% del tiempo). Cada paso tiene su fundamento matemático. A partir de un modelo de la señal
con y sin crisis se deducen los atributos (potencia por banda con la transformada de Fourier y el método de Welch) y una
calibración por paciente, que consiste en restar la mediana de cada atributo. Se desarrollan la optimalidad local y global
y la convexidad, la convergencia del descenso por gradiente, los multiplicadores de Lagrange y las condiciones KKT, el
método simplex y el lema de Neyman-Pearson, que da el umbral de decisión con una tasa de falsas alarmas fija.

En Bonn, con 20 atributos espectrales, la regresión logística alcanza 97,4~\% de exactitud en tres clases y el SVM con núcleo
RBF, 98,4~\%. Se verifica numéricamente que en la logística, que es convexa, 20 puntos iniciales llegan al mismo óptimo, y que
una red neuronal chica termina en 8 mínimos distintos. En CHB-MIT, evaluando sobre un paciente no visto, la logística
calibrada detecta 24 de 34 crisis con 1,1 falsas alarmas por hora.

Una red convolucional sobre espectrogramas se entrena con seis optimizadores (SGD, Momentum, Nesterov, RMSprop, Adam y
AdamW), tres tasas de aprendizaje y cinco semillas. Con la tasa elegida por validación, SGD es significativamente mejor
(AUC de prueba $0{,}976 \pm 0{,}006$). Pero un experimento exploratorio muestra que con la tasa alta SGD, RMSprop, Adam y AdamW
quedan empatados: pesa más la tasa que el algoritmo. El paciente de validación quedó saturado y no sirvió para elegir. En los
pacientes de prueba, la red detecta 10 de 14 crisis con 0,19 falsas alarmas por hora, contra 6 de 14 con 0,54 de la regresión
logística, a pesar de tener menor AUC. Para un detector de crisis, la métrica relevante es la sensibilidad con muy pocas falsas
alarmas, no el AUC.

\medskip
\noindent\textbf{Palabras clave:} epilepsia, EEG, transformada de Fourier, convexidad, multiplicadores de Lagrange,
Neyman-Pearson, redes convolucionales, optimizadores.
